% =====================================================================
\section{Loops and the developmental window}
\fig[0.7\textwidth]{fig_loops.png}{The positive-feedback loops described in reviews, and the place of adolescence on the scale of a lifetime.}
The loop of section~7 is not the only one. Reviews describe several circuits passing through oxidative stress, and all of them converge on the PV cell (fig.~7). A separate question is why the disease begins in youth. The biologists' answer: PV cells and their protective shells mature late, towards the end of adolescence, and are vulnerable until then.
\facts
\begin{itemize}
\item Two principal loops between weak NMDA-receptor function and oxidative stress have been described: an intracellular one (receptor $\to$ protective genes) and a network one (disinhibition $\to$ inflammatory signal $\to$ Nox2 $\to$ oxidants).\tS
\item Another review counts five loops: through the receptor, the mitochondria, the brain's immune cells, the shells around the cells, and dopamine. The review contains no numbers.\tS
\item Genetic risk is usually compensated; the outcome depends on the strength \emph{and the duration} of the stress; mice with part of the receptor knocked out need a ``second hit'' to develop impairments.\tS
\item Oxidative stress damages PV cells specifically in childhood and adolescence.\tS
\item In rats a short course of a sedative in adolescence prevents subsequent impairments.\tS
\end{itemize}
\begin{ineq}
A ``vicious circle'' in a biological review is a set of arrows in a figure. Whether the loop closes depends on the product of the gains along it: if that product is less than one, a perturbation decays. Reviews do not pose this question, because they have no numbers. The developmental window is a slow drift of the defence parameters, that is, a slow passage through the boundary of the region of two states.
\end{ineq}
